\documentclass[openany,fontset=fandol]{ctexbook}
\usepackage{amsmath,amssymb}
\usepackage[unicode,hidelinks]{hyperref}
\setlength{\emergencystretch}{2em}
\interfootnotelinepenalty=10000
\newcommand{\OO}{\mathcal O}
\newcommand{\LL}{\mathcal L}
\newcommand{\MM}{\mathcal M}
\newcommand{\EE}{\mathcal E}
\DeclareMathOperator{\Pic}{Pic}
\DeclareMathOperator{\Sym}{Sym}
\DeclareMathOperator{\Proj}{Proj}
\DeclareMathOperator{\Spec}{Spec}
\DeclareMathOperator{\Hom}{Hom}
\begin{document}
\title{代数几何基础第五卷预备笔记\\射影丛的 Picard 群与零次态射}
\author{据法文扫描和英译编辑稿互校续译}
\date{2026年11月21日}
\maketitle
\chapter*{范围与接续}
本批为第三篇开头：§0 引论及 §1 自定理1.1至零次态射，含中间的回顾、证明和作者改写意见。英稿物理86--90页，法扫描按文本顺序为物理151、152、157、154--156、116--118页；116--118页手写编号123--125。法物理118页只译首行承接上页的句尾，接续从同一行的“2) \(n=1\)”开始。不据扫描的物理顺序重排正文，亦不插入物理153页。

英译编辑者注明下文有一段在手稿中被划去，自 §1 开头至零次态射末止。为保存预备笔记的内容，仍译录这一段，不将其当作后来正式刊定的文本。正文纠正已核实的问题，原读法列入脚注；现代验证与原译分开。本批未并入1158页合订本。
\chapter*{第三篇}
\section*{§0　射影丛与多重射影丛上的可逆层和相对除子；除子的线性系}
\begin{enumerate}
\item 射影丛上的可逆层；应用于从一个射影丛到另一个射影丛的态射。\footnote{此处按法文目录。英稿将本项写成射影与多重射影丛上的可逆层，漏掉应用。}
\item 相对除子函子 \(\operatorname{Div}_{X/S}\) 的可表性；射影丛与多重射影丛上的相对除子。
\item 除子的线性系与到射影丛的态射。
\item 除子的线性系与可逆模。
\end{enumerate}

本节及后面各节展开的结果，已有一部分是整体性的。其中给出了关于射影概形的若干补充：除本第四章中的纯局部结果以外，还用到第二章的整体构造，甚至在这里用到第三章的某些结果。

本节的目标之一是发展“除子的线性系”这一套语言。它一方面关联到向射影丛的态射的分类，另一方面关联到给定相对预概形上可逆模的分类。还应指出，对除子的线性系而言，真正自然的“参数概形”是布劳尔--塞维里概形。这类概形推广了射影丛，例如可定义为这样的纤维丛：经一次满射的平展基变换后，成为与某个射影丛同构的纤维丛。其研究要利用第五章展开的下降理论\footnote{英译括注第五章系指原规划；这不是对后来出版卷次的断定。}，而其分类又等价于以射影线性群为结构群的挠子的分类。因而，我们把这些概形及其与除子的线性系的联系留到本著作第六章研究。

从技术上看，本节的主要结果是定理1.1：用底预概形的 Picard 群确定射影丛的 Picard 群。其最初的一些推论将在第1、2小节给出。

\section*{§1　射影丛上的可逆层与除子；除子的线性系}
\noindent\textbf{1. 确定射影丛上的可逆层；应用于射影丛的自同构。}
以下 \(f:P\to S\) 为结构态射，采用商的约定 \(P=\Proj_S(\Sym \EE)\)；\(\LL(n)=\LL\otimes_{\OO_P}\OO_P(n)\)。

\medskip\noindent\textbf{定理1.1。}
设 \(S\) 是预概形，\(\EE\) 是 \(S\) 上的局部自由模，在每点处的秩均有限且至少为 \(2\)。令 \(P=\mathbf P(\EE)\) 为它所定义的射影丛。对 \(P\) 上任一可逆模 \(\LL\)，存在一族由 \(\mathbb Z\) 编号的两两不交开集 \((S_n)_{n\in\mathbb Z}\)，覆盖 \(S\)，以及 \(S\) 上的可逆模 \(\MM\)，使得在各 \(f^{-1}(S_n)\) 上有
\[
\LL|_{f^{-1}(S_n)}
  \simeq \bigl(f^*\MM\otimes_{\OO_P}\OO_P(n)\bigr)|_{f^{-1}(S_n)} .
\]
这族 \(S_n\) 由上述条件唯一确定；连同所给分层同构，\(\MM\) 也在唯一的相容同构意义下确定。\footnote{法扫描157页与英稿87页的扭转印作 \(\OO_P(1)\)，但同一式右侧写 \(f^*(\MM)(n)\)，且后证用消去扭转 \(\LL(-n)\)。故应为 \(\OO_P(n)\)。原稿“在唯一同构意义下”须连同所给同构理解；未加相容条件的线丛本身一般有非平凡自同构。}

\medskip\noindent\textbf{评注1.2。}
若不对 \(\EE\) 的秩作假设，\(S\) 典范分解为三个开子预概形的不交和 \(S^0,S^1,S^2\)，使 \(\EE\) 在这三部分上的秩分别为 \(0,1,\ge 2\)。于是，确定 \(P\) 上的可逆模归结为分别确定 \(f^{-1}(S^i)\)（\(i=0,1,2\)）上的可逆模。\(i=2\) 可用1.1；\(f^{-1}(S^1)\) 在 \(S^1\) 上同构于 \(S^1\)，其 Picard 群就是 \(\Pic(S^1)\)；最后 \(f^{-1}(S^0)\) 为空，其 Picard 群为零。

\medskip\noindent\textbf{推论1.3。}
在1.1的假设下，再设 \(S\) 连通且非空。则 \(P\) 上的每个可逆模 \(\LL\) 都同构于 \(f^*(\MM)(n)\)，其中 \(n\in\mathbb Z\)，\(\MM\) 是 \(S\) 上的可逆模。\(n\) 唯一，而 \(\MM\) 连同这个分解在唯一的相容同构意义下确定。

也可以这样表述该推论。考虑自然同态 \(\Pic(S)\to\Pic(P)\) 和 \(\mathbb Z\to\Pic(P)\)：前者来自 \(f:P\to S\)，后者由 \([\OO_P(1)]\) 确定。于是有典范同态
\[
(*)\qquad
\Pic(S)\times\mathbb Z\longrightarrow\Pic(P),\qquad
([\MM],n)\longmapsto[f^*(\MM)(n)] .
\]
定义它并不要求 \(S\) 或 \(\EE\) 满足额外条件。

\medskip\noindent\textbf{推论1.4。}
在1.1的条件下，若 \(S\ne\varnothing\)，则上述同态是单射；若 \(S\) 还连通，则为双射。

若去掉秩至少为 \(2\) 的假设，由1.2和1.4可知，\(S\) 连通时该同态仍满射，却不一定单射：当 \(\EE\) 的秩为 \(1\) 时，其核同构于 \(\mathbb Z\)；秩为 \(0\) 时，核就是 \(\Pic(S)\times\mathbb Z\)。

\medskip\noindent\textbf{定理1.1的证明。}
先证唯一性。若 \(S\) 是域 \(k\) 的谱，只有 \(n=m\) 时 \(\OO_P(n)\) 才可能同构于 \(\OO_P(m)\)，亦即只有 \(m-n=0\) 时 \(\OO_P(m-n)\) 才可能同构于 \(\OO_P\)。这是因为 \(\OO_P(1)\) 丰沛，而且由秩假设有 \(\dim P\ge 1\)。交换 \(m,n\) 后可设 \(d=m-n\ge 0\)。若 \(d>0\)，则 \(\OO_P(d)\) 丰沛；它若同构于 \(\OO_P\)，\(P\) 就必须拟仿射，又因为 \(P\) 在 \(k\) 上固有，便是有限 \(k\)-概形，与 \(\dim P\ge 1\) 矛盾。因此，把分层同构限制到纤维就能证出 \((S_n)\) 的唯一性。

至于 \(\MM\) 的唯一性，可归结到 \(S=S_n\)。此时给定 \(\LL\simeq f^*(\MM)(n)\)，有由此唯一确定的同构
\[
(**)\qquad \MM\xrightarrow{\sim}f_*(\LL(-n)).
\]
事实上，消去扭转得 \(\LL(-n)\simeq f^*\MM\)，推前后把右项与 \(f_*f^*\MM\) 等同；由于 \(\MM\) 局部自由，投影公式给出 \(f_*f^*\MM\simeq \MM\otimes f_*\OO_P\)，而 \(\OO_S\xrightarrow{\sim}f_*\OO_P\)。取相应方向即得 \((**)\)。

下面证明 \((S_n)\) 与 \(\MM\) 存在。借助已经证明的唯一性，这是 \(S\) 上的局部问题，因而归结为：

\medskip\noindent\textbf{推论1.5。}
设 \(\EE\) 是 \(S\) 上的有限秩局部自由模，\(P=\mathbf P(\EE)\)，\(\LL\) 是 \(P\) 上的可逆模，且 \(s\in S\)。那么存在 \(s\) 的开邻域 \(U\) 和整数 \(n\)，使
\[
 \LL|_{f^{-1}(U)}\simeq \OO_P(n)|_{f^{-1}(U)} .
\]
当 \(\EE\) 在 \(s\) 的秩至少为 \(2\) 时，\(n\) 唯一；当秩至多为 \(1\) 时，推论本身局部立即成立。\footnote{法稿把秩至多为 \(1\) 写作另一情形，不是前一句的附加假设。秩 \(0\) 时射影丛为空；秩 \(1\) 时 \(f\) 局部是 \(U\) 本身，\(\LL\) 与 \(\OO_P(1)\) 缩小 \(U\) 后均平凡。}

由于问题是局部的，可设 \(\EE=\OO_S^{r+1}\)，从而 \(P=\mathbf P_S^r\)。再用第8节的有限逼近办法，可归结到 \(S\) 诺特的情形。证明分两步。\footnote{原稿简指第8节，英编辑稿对卷次附问。此处沿用手稿的节号；校核附记给出这里所需的有限表示下降。}

\medskip\noindent\textbf{a) \(S\) 是域 \(k\) 的谱。}
记 \(A=k[t_0,\ldots,t_r]\)，赋予通常的分次，记 \(V=\mathbb A_k^{r+1}\setminus\{0\}\)，\(i:V\hookrightarrow\mathbb A_k^{r+1}\) 为开浸入，\(q:V\to\mathbf P_k^r\) 为典范投影。我们知道，\(\LL\) 可由 \(A\) 上有限生成的分次模 \(L\) 给出；\(L\) 对应的仿射拟凝聚层在 \(V\) 上的限制 \(\widetilde L'\)，恰为 \(q^*\LL\)，因而可逆。\(A\) 是唯一分解整环，其在原点的局部化仍是唯一分解整环，且维数至少为 \(2\)。由此 \(i_*\widetilde L'\) 是可逆模，相应于 \(A\) 上的可逆模
\[
N=\Gamma(V,q^*\LL).
\]
\(N\) 自然带有分次；自然映射 \(L\to N\) 在仿射空间中原点以外的各点上都是同构。因此，用 \(N\) 取代 \(L\)，可归结到 \(L\) 可逆的情形。由于 \(A\) 是唯一分解整环，先不计分次，\(L\) 是秩一自由模。一个标准引理（想必可在布尔巴基著作中找到）表明它作为分次模也是秩一自由的，故在 \(P\) 上相伴的模同构于某个 \(\OO_P(n)\)。

\emph{作者的行文意见。}从写作上看，从一个有限生成的 \(L\) 开始并不妥当。不如直接从 \(L=\Gamma(V,q^*\LL)\) 开始，它给出 \(i_*q^*\LL\)；第二章表明这个分次模恰好确定 \(\LL\)。随后用上述论证证明它是秩一分次自由模即可。\footnote{法稿这一段给 \(i_*\) 处写作 \(q_*\)，英稿也有推、拉混写；这里的层在仿射空间上，故应为开浸入 \(i\) 的推前。作者所说布尔巴基的具体条目没有核定。}

\medskip\noindent\textbf{b) 一般情形。}
由 a) 和 III.4.6.5 可得，所用的关系是第三章第2节所证的
\(\mathrm H^1(\mathbf P_k^r,\OO_{\mathbf P_k^r})=0\)。
\hfill 证毕。

\medskip\noindent\textbf{1.4 的另一种表述。}
令 \(\mathbb Z(S)\) 表示 \(S\) 上的局部常值整数函数群。其元素 \(\nu\) 对应于两两不交的开分割 \((S_n)\)（允许部分为空）；在 \(f^{-1}(S_n)\) 上取 \(\OO_P(n)\) 就定义了可逆模 \(\OO_P(\nu)\)。这给出自然同态 \(\mathbb Z(S)\to\Pic(P)\)，并给出变式
\[
(*\mathrm{\ bis})\qquad \Pic(S)\times\mathbb Z(S)\longrightarrow\Pic(P).
\]
更一般也更令人满意的断言是：在1.1的假设下，此同态为双射。附记：把整数送到对应常值函数的典范映射 \(\mathbb Z\to\mathbb Z(S)\)，在 \(S\) 非空时为单射，在 \(S\) 连通时为满射；于是形式上恢复1.4，宜将1.4写成上述较一般的形式。\footnote{法稿156页手写补入“分别，连通”，英稿漏录连通条件。变式的记号按法稿为 \(\mathbb Z(S)\)，不取英稿的 \(\mathbb Z(A)\)。}
还应评注：用第五章1.1将引入的 Picard 概形语言，该等价表述就是说，由 \(\OO_P(1)\) 所定义的截面诱导的典范同态
\(\mathbb Z_S\to\Pic_{P/S}\)
是同构；这里 \(\mathbb Z_S\) 是 \(S\) 上由 \(\mathbb Z\) 给出的常值群概形。

\medskip\noindent\textbf{可逆模的度数。}
设 \(P\) 是域 \(k\) 上的射影丛。若 \(P\) 上的可逆模 \(\LL\) 同构于 \(\OO_P(n)\)，便称它为 \(n\) 次。当 \(\dim P\ge 1\) 时，\(\LL\) 确定 \(n\)；若 \(\dim P\le 0\)，即 \(P\) 为空或只是一点，则 \(\LL\) 对每个整数 \(n\) 都是 \(n\) 次。

若 \(P\) 是任意底概形上的射影丛，称 \(\LL\) 为 \(n\) 次，是指它在每个纤维上均为 \(n\) 次。除非 \(P/S\) 的所有纤维维数都至多为 \(0\)，否则 \(n\) 由 \(\LL\) 确定；在这一例外情形下，\(\LL\) 对任何整数 \(n\) 都是 \(n\) 次。\footnote{按法扫描116页译出关于任意底的定义及唯一性，英稿89页跳过了这几行。}
由1.1与1.2，\(\LL\) 为 \(n\) 次也就是说，\([\LL]\) 在 \((*)\) 下属于 \(\Pic(S)\times\{n\}\) 的像，即同构于 \(f^*(\MM)(n)\)。若所有纤维非空，即 \(\EE\) 处处秩至少为 \(1\)，则对固定的 \(n\)，\(\MM\) 连同相容同构由 \(\LL\) 唯一确定；这仍由1.1和1.2得到。

顺便说，定理1.1其实宜不加秩限制地陈述：\(P\) 上每个可逆模 \(\LL\) 都可写成原陈述的形式；若 \(P\) 的纤维均非空，亦即 \(\EE\) 处处非零秩，则选定扭转后，\(\MM\) 连同相容同构由 \(\LL\) 确定；若 \(P\) 的全部纤维维数至少为 \(1\)，亦即 \(\EE\) 的秩处处至少为 \(2\)，则底的分割也由 \(\LL\) 唯一确定。这样，评注1.2可移入证明。\footnote{法稿和英稿这段均误写“\(S\) 上的 \(\LL\)”；应在 \(P\) 上。法稿还分别给出两种秩条件，英稿90页合并成“非空即秩至少为2”，混淆了基底可逆模与分割的唯一性。秩1时不同扭转会改变 \(\MM\)，故必须固定扭转。}

\medskip\noindent\textbf{射影丛之间的态射；零次情形。}
设 \(P'=\mathbf P(\EE')\) 是第二个射影丛。确定 \(\Pic(P)\) 原则上即可确定 \(S\)-态射 \(g:P\to P'\)。这样的态射由 \(P\) 上的可逆模 \(\LL\simeq g^*\OO_{P'}(1)\) 及同态 \(\EE'\to f_*\LL\) 确定，要求伴随同态 \(f^*\EE'\to\LL\) 满射，并把相差一个 \(\LL\) 的相容同构的表示视为相同。若 \(\LL=g^*\OO_{P'}(1)\) 为 \(n\) 次，则称 \(g\) 为 \(n\) 次。于是只需对各给定 \(n\) 分类这些态射。若 \(P\) 的纤维维数皆至少为 \(1\)，则 \(g\) 为 \(n\) 次必有 \(n\ge0\)，因为域上正维射影空间中的 \(\OO(n)\) 只有在 \(n\ge0\) 时才由其截面生成。自然也可以按1.2分别讨论各秩层，以下只考虑纤维维数至少为 \(1\) 的情形。

由于
\[
 f_*\bigl(f^*(\MM)(n)\bigr)
   \simeq \MM\otimes f_*\OO_P(n)
   \simeq \MM\otimes\Sym^n(\EE),
\]
确定 \(g\) 归结为确定相差同构的一对 \((\MM,u)\)：\(\MM\) 是 \(S\) 上的可逆模，\(u:\EE'\to\MM\otimes\Sym^n(\EE)\) 是同态，且对应的 \(f^*\EE'\to f^*(\MM)(n)\) 为满射。\footnote{英稿此处将推前的最后一项印作 \(\Sym^2(\EE)\)；法稿117页和随后给出的 \(u\) 都是 \(\Sym^n(\EE)\)。目标的扭转应是 \(\OO_{P'}(1)\)，不是 \(\OO_P(1)\)。}
所以，\(g\) 确定一个整体不变量 \([\MM]\in\Pic(S)\)。给定一个代表 \(\MM\) 后，相应的 \(g\) 由下述商集的某个子集给出：
\[
 \Hom_{\OO_S}(\EE',\MM\otimes\Sym^n(\EE))
   \big/\Gamma(S,\OO_S)^\times .
\]
取这个商，正对应于用可逆模描述 \(g\) 时的“模同构”。附记：\(P\) 上可逆模的自同态及自同构，分别对应于 \(\OO_P\) 的整体截面及可逆整体截面；因为 \(\OO_S\simeq f_*\OO_P\)，也分别对应于 \(\OO_S\) 的整体截面及可逆整体截面。

\medskip\noindent\textbf{特殊情形 1)：\(n=0\)。}
这时所取的是满射 \(\EE'\twoheadrightarrow\MM\)，即在各点的剩余域上非零的同态，模去 \(\MM\) 的同构。我们所得的恰是 \(g=h\circ f\)，其中 \(h:S\to P'\) 为截面，即由 \(\EE'\) 的可逆商 \(\MM\) 所确定的截面。故从 \(P\) 到 \(P'\) 的零次 \(S\)-态射恰为相对于 \(S\) 的常值态射。\footnote{法稿117页把“处处非零”改成满射，意为在每个剩余域上的映射非零，不是只看茎上的映射是否为零。法稿的 \(hf\) 与 \(P'\) 修复了英稿的 \(h\) 与 \(P\)。该结论沿用前段约定；在正秩层上，\(f\) 忠实平坦，能检验满射。}

\emph{〔英译编辑者注明：手稿划去的段落到此为止。〕}

\section*{校核附记（非手稿）}
\paragraph{分层与相容同构。}
由1.5知，在秩至少为2的部分，纤维次数局部常值。对每个 \(S_n\) 令 \(P_n=f^{-1}(S_n)\)、\(f_n=f|_{P_n}\)，取
\[
\MM_n=(f_n)_*\bigl(\LL|_{P_n}(-n)\bigr).
\]
由1.5在 \(S_n\) 上的局部平凡性，\(\MM_n\) 可逆；伴随评价映射在这些开集上是同构，所以全局给出 \(f_n^*\MM_n(n)\simeq\LL|_{P_n}\)。这些 \(\MM_n\) 在不交开分割上拼成 \(\MM\)。若只给一个整数 \(n\)，在任一非空、处处正秩的底上，使用同一个 \(f_*(\LL(-n))\) 即可；不能把这个说法解释为可逆模没有标量自同构。底为两部分不交并、两部分上的次数分别为0和1，就是必须使用局部常值次数的例子。

\paragraph{秩为零或一。}
在商约定下，秩一时典范满射 \(f^*\EE\to\OO_P(1)\) 是秩一模之间的同构。识别 \(P=S\) 后，\((*)\) 就是
\(([\MM],n)\mapsto[\MM\otimes\EE^{\otimes n}]\)，核为
\(\{([\EE^{\otimes(-n)}],n):n\in\mathbb Z\}\)。
秩零时 \(P\) 为空，故核是整个源。处处至少一秩时，\(f_*\OO_P\simeq\OO_S\)； 秩为零时不满足此式（除非底也为空）。

\paragraph{域上与局部提升的两个细节。}
当 \(r\ge1\)，删去原点没有删去任何余维一点；在正规仿射空间上，余维一的除子数据决定秩一反身层。由 \(A\) 唯一分解，\(V\) 的任何 Cartier 除子都主，故 \(q^*\LL\) 无分次地平凡；又 \(\Gamma(V,\OO_V)=A\)（在分式域内取余维一局部化之交），所以 \(i_*q^*\LL\) 平凡。若分次 \(A\)-模 \(N\) 无分次地自由秩一，\(N/(t_0,\ldots,t_r)N\) 是一维分次 \(k\)-向量空间。提升它的齐次基元；由于 \(N\) 有限生成且分次下有界，分次 Nakayama 引理说明所提升的元生成 \(N\)，于是它是齐次基。这样的移位恰给 \(\OO(n)\)。

对一般的底，在仿射邻域上可把一个可逆模及其乘法逆的给定同构从 \(\mathbf P_A^r\) 沿有限标准仿射覆盖下降到某个有限生成 \(\mathbb Z\)-子代数：有限表示、有限重叠上的过渡映射及有限等式都只需有限多个系数。诺特情形下，对 \(s\) 的纤维所给的平凡化，在扭掉 \(n\) 后由上同调与基变换及 \(H^1=0\) 可提升为邻域上的截面。它的非生成点构成闭集；因 \(f\) 固有，其像闭且不含 \(s\)。移去该像，截面在整个限制空间上给出平凡化。这说明1.5中“纤维平凡”何以能加强成“在底的邻域上的整个限制空间平凡”。

\paragraph{底本与引用。}
法稿据随工程所附按原文顺序摘出的扫描；英稿见
\href{https://webusers.imj-prg.fr/~leila.schneps/grothendieckcircle/EGA/EGA5trans.pdf#page=86}{第五卷英译编辑稿物理86页起}。
射影丛的商约定和 \(\OO(1)\) 见
\href{https://stacks.math.columbia.edu/tag/01OA}{Stacks Project, Tag 01OA}。
手稿中的 III.4.6.5 与第五章1.1均保留原写法，不把现行编号移入原引文。
\end{document}
